\section{Discussion and limitations}
\label{sec:discussion}

\paragraph{What $\Rmodel$ adds.}
The metric is most useful as a common accounting layer.  It makes the workload,
operation set, causal mask, transformed operands, and dense head explicit; it
then permits sitewise and operationwise diagnoses below the scalar summary.
The \Afour/\Aseven{} comparison shows why this hierarchy matters: local zero
percentages identify where representations collapse, while per-operation counts
identify why the graph-level opportunity changes.  A single number without
those decompositions would hide the mechanism; local plots without the common
denominator would not reveal its computational footprint.

\paragraph{How we use 410M.}
Excluding the cohort would produce a cleaner but misleading scale-persistence
narrative.  Including it as ordinary replication would overstate the design.
We instead use it as a stress test: the high-dose topological advantage
persists, but the trained loss response reverses under sharply lower exposure.
The disagreement with post-hoc clipping suggests that learned adaptation to a
gate can matter more than instantaneous deletion.  It does not distinguish
beneficial regularization, altered optimization geometry, capacity effects, or
undertraining.  A matched tokens-per-parameter or loss-matched study and
replicate seeds are required before any scaling statement.

\paragraph{Sitewise redistribution is selective.}
The cross-scale endpoint heatmap shows that broader pressure does not induce a
uniform attention response.  An earlier matched 14M diagnostic reinforces this
point: as near-zero mass at directly pressured $h$ rose from 63.5\% to 92.4\%,
$v$ moved only from .133\% to .185\%; $q$, $k$, and $m$ remained tiny and
non-monotone, while the post-$W_o$ attention output rose through the middle
doses and then fell.  We treat ``spillover'' as this site-specific descriptive
redistribution, not a causal propagation mechanism (Appendix
Figure~\ref{fig:spillover}).

\paragraph{Limits of the intervention evidence.}
Each endpoint has one seed.  \Afour{} and \Aseven{} differ jointly in gates and
pressure sites.  The pressure weight and trust budget were not broadly tuned.
A matched 14M comparison of naive and orthogonal $\ell_1$ pressure produced
small, mixed loss differences, so the present data do not establish OL1 as
superior (Appendix~\ref{app:ol1}).  Post-hoc thresholds are calibrated on ten
training blocks, not fit to validation, but they remain a particular clipping
rule rather than a universal comparator.

\paragraph{Limits of the computational claim.}
Our denominator describes complete length-2,048 sequences with uncached causal
attention.  It resembles prefill or training accounting, not cached
autoregressive decoding, whose KV reuse changes operation footprints.  Exact
zeros alone do not create sparse execution.  We report no removed FLOPs,
latency, energy, memory-traffic benefit, or end-to-end speedup.  Establishing
those outcomes requires kernels whose format and scheduling match the measured
site structure.

